% Preamble
\documentclass[twocolumn]{article}

% Packages
\usepackage[T1]{fontenc}
\usepackage[utf8]{inputenc}
\usepackage{indentfirst}
\usepackage{amsmath}
\usepackage[margin=1in]{geometry}
\usepackage[font=small,labelfont=bf]{caption}
\usepackage{hyperref}
\usepackage{fancyhdr}
\usepackage{booktabs}
\usepackage{graphicx}
\usepackage{listings}
\usepackage{xcolor}

\lstset{
    basicstyle=\ttfamily\footnotesize,
    breaklines=true,
    columns=fullflexible,
    keepspaces=true,
    language=C++,
    showstringspaces=false,
    frame=single,
    xleftmargin=0.5em,
}

\pagestyle{fancy}

\fancyhead{}
\lhead{
    \includegraphics[scale=0.13]{images/logo.png}
}
\rhead{
    Distributed and Parallel Computing \\
    2019/2020
    \footnotesize
}

\fancyfoot[C]{\thepage}

\date{}

\title{
    \includegraphics[scale=0.2]{images/logo.png}\\
    \vskip 0.5em
    \Huge Distributed and Parallel Computing\\
    \LARGE MPI + OpenMP Hybrid Report\\
}

\author{
    Pedro Rio\\
    \textit{MSc Data Science and Engineering}\\
    \textit{Instituto Superior Técnico}\\
    \textit{97241}
    \and
    Sebastião Maia Serqueira\\
    \textit{MSc Data Science and Engineering}\\
    \textit{Instituto Superior Técnico}\\
    \textit{97108}\\
    \and
    Eduardo Machado\\
    \textit{MSc Data Science and Engineering}\\
    \textit{Instituto Superior Técnico}\\
    \textit{97133}
}

% Document
\begin{document}
    \maketitle


    \section{Program Overview}\label{sec:program-overview}
    The problem is a matrix-factorisation recommender: a sparse user--item matrix $A$ is
    approximated by $A \approx L \cdot R$, with $L$ of size
    $\textit{users} \times \textit{features}$ and $R$ of size
    $\textit{features} \times \textit{items}$. The factors are fitted by gradient descent
    over the observed entries, and the final dense product $B = L \cdot R$ yields one
    recommendation per user. The four stages are \emph{ReadInput}, \emph{InitialLR},
    \emph{UpdateLR} (repeated for the requested number of iterations), and
    \emph{FilterFinalMatrix}.

    This report covers the \textbf{MPI + OpenMP hybrid} deliverable. Its distribution
    across processes --- the 2D grid, the partitioning and the row/column reductions
    described in the next two sections --- is shared with the pure MPI deliverable in
    \texttt{mpi/}, which has its own report. What is specific to this variant is the
    shared-memory layer inside each rank, in Section~\ref{sec:hybrid}, and the thread
    behaviour measured in Section~\ref{sec:evaluation}.


    \section{Why a 2D Grid}\label{sec:why-grid}
    The simplest MPI design is \emph{replicate-all}: every process keeps full copies of
    $L$ and $R$, splits the non-zero list into blocks, computes a partial gradient, and
    sums it with a global \texttt{MPI\_Allreduce}. It is correct and short, but it does not
    scale: the reduced gradient has $O(\textit{features} \cdot (\textit{users} +
    \textit{items}))$ entries and that volume \emph{does not shrink} as processes are
    added, while compute per process falls as $1/P$. Communication therefore dominates at
    scale, and on \emph{wide} instances (a large $R$) it \emph{anti-scales} --- adding
    processes makes it slower. Memory does not scale either: a single node must hold the
    whole problem.

    Both limits have the same cause --- replication --- so we distribute the factors over a
    2D process grid.


    \section{The 2D-Grid Version}\label{sec:grid}

    \subsection{Process mesh}\label{subsec:mesh}
    \texttt{MPI\_Dims\_create} factors $P$ into a balanced $P_r \times P_c$ grid;
    \texttt{MPI\_Cart\_create} builds the mesh (with reorder disabled, so a rank's grid id
    equals its world rank) and \texttt{MPI\_Cart\_sub} derives a \emph{row} and a
    \emph{column} sub-communicator for each rank.

    \subsection{Foster methodology}\label{subsec:foster}
    \textbf{Partitioning.} $L$ is split by \emph{user-rows} across the grid rows and $R$ by
    \emph{item-columns} across the grid columns. A non-zero $(u,i)$ is owned by rank
    $(\textit{userBlock}(u), \textit{itemBlock}(i))$, so each rank holds exactly the
    $L$-rows and $R$-columns its non-zeros touch --- memory per rank is
    $O((\textit{users} \cdot \textit{features} + \textit{features} \cdot
    \textit{items})/\sqrt{P})$.

    \textbf{Communication.} Each rank computes a partial gradient from its cell's
    non-zeros (all data local), then $dL$ is summed across the \emph{row} communicator and
    $dR$ down the \emph{column} communicator. These are reductions of $\sqrt{P}$-smaller
    slices over $\sqrt{P}$-sized sub-communicators, not a global reduce of the whole
    gradient.

    \textbf{Agglomeration and Mapping.} The non-zeros are agglomerated into one contiguous
    block per grid cell and routed there once with a bucketed \texttt{MPI\_Scatterv}; one
    cell maps to one process.

    \subsection{Deterministic, communication-free init}\label{subsec:init}
    \texttt{initialLR} is seeded with a fixed value, but a partitioned factor cannot simply
    seek into the random stream (an $R$-column slice is strided). The trick is
    \emph{generate-and-discard}: every rank runs the same \texttt{srandom(1)} sequence in
    the serial fill order but stores only the values that fall in its slice. The slices
    hold exactly the serial values, with no broadcast of the factors.

    \subsection{Distributed final stage}\label{subsec:final}
    $B = L \cdot R$, the mask of known entries, and the per-user arg-max are all computed
    on the cells: each rank forms its $B$-block, takes a local arg-max over its item
    columns, and the results are combined across the row with \texttt{MPI\_MAXLOC}
    (smallest item index on ties, matching the serial scan) and gathered to the root in
    user order. There is no serial-root tail.


    \section{Hybrid OpenMP Layer}\label{sec:hybrid}
    Each rank can fan its local work across OpenMP threads. A \emph{single} parallel region
    is hoisted around the whole iteration loop (the team forks once), \texttt{updateLR}
    uses orphaned worksharing, and the per-thread gradient conflict is resolved with an
    array reduction --- \texttt{reduction(+ : dL[:lN], dR[:rN])} --- exactly as in the
    shared-memory version. MPI stays on the main thread under
    \texttt{MPI\_THREAD\_FUNNELED}: the two \texttt{Allreduce} calls run inside an
    \texttt{omp master} block, fenced by barriers. The default is one thread per rank (pure
    MPI); threads are explicit opt-in through \texttt{OMP\_NUM\_THREADS}, and one thread
    reproduces the pure-MPI behaviour exactly.


    \section{Evaluation}\label{sec:evaluation}
    The question this deliverable has to answer is not whether it is faster than serial
    --- it is --- but whether putting threads \emph{inside} the ranks beats spending the
    same cores on more ranks. That is only answerable at equal core counts, so
    Table~\ref{tab:hybrid-vs-pure} compares the hybrid at two threads per rank against the
    pure MPI deliverable at the same total core count.

    \textbf{Provenance.} Both designs were measured in the same run by
    \texttt{compare/scripts/benchmark.py} on an Apple Silicon laptop (macOS 15.6, 10
    cores, 8 of them performance cores), three repetitions with the minimum kept, every
    run verified against the instance's \texttt{.out} file. The figures are the committed
    \texttt{compare/data/comparison.hybrid.csv} and
    \texttt{comparison.mpi.csv} and are reproducible from this repository.

    \begin{table}[!h]
        \caption{\label{tab:hybrid-vs-pure}Hybrid ($2$ threads/rank) vs.\ pure MPI at
        equal core count (time, s). Both measured on this laptop in one run.}
        \centering
        \resizebox{\linewidth}{!}{
            \begin{tabular}[t]{lrrrrrr}
                \toprule
                & \multicolumn{2}{c}{2 cores} & \multicolumn{2}{c}{4 cores}
                & \multicolumn{2}{c}{8 cores}\\
                \cmidrule(lr){2-3}\cmidrule(lr){4-5}\cmidrule(lr){6-7}
                instance & hybrid & pure & hybrid & pure & hybrid & pure\\
                \midrule
                inst200 (wide) & 4.29 & \textbf{3.94} & 4.90 & \textbf{2.25} & \textbf{2.90} & 2.94\\
                inst500        & 5.55 & \textbf{4.76} & 3.72 & \textbf{2.66} & 3.84 & \textbf{3.30}\\
                inst1000       & 2.45 & \textbf{2.22} & 1.66 & \textbf{1.23} & 1.47 & \textbf{0.99}\\
                \bottomrule
            \end{tabular}}
    \end{table}

    \textbf{On a single node, pure MPI wins almost everywhere.} The hybrid is $9$--$17\,\%$
    slower at two cores, $35$--$118\,\%$ slower at four, and $16$--$48\,\%$ slower at
    eight; the single exception is the wide instance at eight cores, where the two are
    level ($2.90$ against $2.94$\,s, a $1.6\,\%$ difference that is inside this machine's
    run-to-run spread). That is the expected result and worth stating plainly: within one
    node MPI already communicates through shared memory, so a rank boundary costs almost
    nothing, while a thread boundary costs a barrier and a reduction on every iteration.

    The one place threads nearly pay for themselves is also the informative one. The wide
    instance ($200$ users, $10{,}000$ items) has the most lopsided grid, so adding ranks
    adds communication faster than it adds useful work; there, trading ranks for threads
    stops being a loss. That is the shape of problem --- and, on a cluster, the network
    cost --- that the hybrid exists for.

    \textbf{What this measurement cannot show.} Every figure here is from one shared-memory
    node, which is precisely the setting where the hybrid has least to offer. Its case is
    one rank per node or socket with threads within, so that inter-node
    \texttt{Allreduce}s run over a $\sqrt{P}$-smaller grid and intra-node work never
    touches the network. Nothing in this repository can measure that, and no figure above
    should be read as evidence about it either way.

    \section{Implications for Parallel and Distributed Computing}\label{sec:pdc}

    \subsection{Distribute the data, not just the work}\label{subsec:distribute}
    Replicate-all distributes the \emph{work} but not the \emph{data}, so it hits both a
    memory wall and a communication ceiling. The 2D grid distributes the data: the factor
    slices and the gradient fall as $1/\sqrt{P}$, and the dense $B$ built in the final
    stage --- which dominates the footprint on every instance here --- falls as $1/P$. In
    this variant there is one more term the pure version does not pay: the gradient loop
    reduces into per-thread private copies, so each thread adds another
    $(U\!\cdot\!F + F\!\cdot\!I)/\sqrt{P}$ to its rank. Threads buy parallelism with
    memory, which is worth knowing before raising the thread count on a tight node.

    \subsection{Match the collective to the decomposition}\label{subsec:collective}
    The win is not a faster \texttt{Allreduce}; it is a \emph{smaller} one, over a
    \emph{smaller} communicator. Because a user's updates are spread only across its grid
    row and an item's only down its grid column, the global reduction decomposes into
    independent row and column reductions. Choosing the communicator to match the data
    layout is what turns flat communication into communication that scales.

    \subsection{Remove the serial tail}\label{subsec:tail}
    Computing $B$ and the arg-max on the root would be an
    $O(\textit{users} \cdot \textit{items} \cdot \textit{features})$ sequential tail that,
    by Amdahl's law, caps the speedup. Distributing it over the same grid keeps the whole
    pipeline parallel.

    \subsection{The hybrid is for the memory hierarchy}\label{subsec:hierarchy}
    One rank per node (or socket) $\times$ threads-within keeps intra-node work in shared
    memory and shrinks the process grid, so the inter-node \texttt{Allreduce}s are cheaper.
    On a single node, where MPI already uses shared-memory transport, the hybrid is a
    wash; its value appears across the network, which is exactly where the cost model says
    fewer, fatter ranks help.

    \subsection{The thread count cannot be guessed at runtime}\label{subsec:threads}
    The best \texttt{OMP\_NUM\_THREADS} depends on memory bandwidth, NUMA, the
    interconnect, and the problem shape --- machine-specific quantities a single run cannot
    know. (A formula $T \approx \sqrt{\textit{nnz}/(\textit{uLocal}+\textit{iLocal})}$ was
    measured $20$--$53\,\%$ off the best in half the cases.) The program therefore defaults
    to one thread and treats threads as explicit configuration: set the rank/thread split
    in the job script and confirm it with a short scaling sweep.

    \subsection{Reproducibility under reordering}\label{subsec:repro}
    The partitioned reductions sum the gradient in a different order than the serial code,
    so floating-point results differ in the last bits --- but the arg-max is robust to
    that, so the recommendations match the serial output bit-for-bit at every process
    count, grid shape, and thread count. A deterministic serial reference is what makes
    this verifiable.

    The guarantee is about the \emph{output}, not the arithmetic. With threads on, $L$ and
    $R$ differ in their last bits between runs: OpenMP does not specify the order in which
    per-thread reduction privates are merged, and implementations merge them as threads
    finish. Measured here, a fixed configuration produced eight distinct $L$/$R$ checksums
    in eight runs --- under \texttt{schedule(static)} and \texttt{schedule(guided)} alike,
    so the loop schedule does not control it. The recommendations are integer arg-max
    indices and are insensitive to drift that small, which is why the byte-for-byte
    guarantee survives. A bitwise-reproducible gradient would need a fixed-order
    reduction.

\end{document}
